\section{Evaluation Methodology}
\label{sec:design}
\subsection{Comparison and Evidence Design}
The principal comparison is between Jev and DeepSeek configured for concise
structured output without a generated reasoning trace. Both use the same
intent contract, admission limits, execution policy, and caching mechanism.
The real-service comparison adds a self-hosted Qwen model and a fixed rule
parser. Model identifiers, software,
hardware, workload parameters, and execution procedures are provided in
Appendix~\ref{app:settings}.

Study A measures live API decisions and
combines their actual timing with modeled service execution. It evaluates
the interpretation tradeoff and how decision waiting consumes execution
slack. Study B executes a real two-node OCR service and checks the returned
text. It tests whether lower decision latency remains useful after actual
communication and processing. The measured hosted latency includes provider
execution, routing, and network transport in both studies.

\subsection{Workloads and Matched Comparisons}
Study A evaluates 216 synthetic English requests covering the full intent
contract. The same requests recur across three consecutive measurement
blocks, with input order varied and backend call order alternated within
pairs. Thus the blocks probe within-session variation on a shared corpus.
An agent generates both the descriptions and their reference labels.

Service conditions vary three factors: steady or bursty arrivals, changing
or repeated descriptions, and caching disabled or enabled. Each of the eight
conditions has a matched workload for the two hosted backends. In Study A,
live admission events feed modeled execution on the same elapsed-time axis;
no images are processed in that study. Study B holds the corresponding
request descriptions, image assignments, arrival schedules, and service
configuration fixed across all four backends. Complete backend runs are
executed sequentially in randomized order. Consequently, matching controls
the workload but does not eliminate variation in provider or network state.

The real service uses Tesseract recognition~\cite{smith2007overview} on images
from the IIIT5K scene-text dataset~\cite{mishra2012scene}. Separate development
images calibrate the scheduler's service-duration estimates. Evaluation
images are selected before recognition results are observed and are not
filtered by whether a tier recognizes them correctly. The tiers are execution
configurations, rather than guarantees of recognition quality. Each condition
includes 36 supported OCR requests and 12 unsupported requests. Images and
intent descriptions are reused across conditions, while a new service run
is still required after an interpretation-cache hit. Deadlines and offered
loads are controlled study parameters.

\subsection{Correctness and Completion}
\label{sec:metrics}
Exact semantic correctness requires all four interpreted fields to equal
their reference values. Study A measures it on the complete paired request
set. Operational completion instead requires a supported request to finish
modeled execution by its deadline while satisfying the reference service,
locality, minimum tier, and mapped priority. Its strict companion additionally
requires exact field equality. Study B's correct completion adds equality
between the returned OCR text and the reference text after canonical Unicode
normalization and removal of surrounding whitespace. Case and internal
characters are retained. This distinguishes interpretation
accuracy, execution compliance, and service-output correctness.

For a given backend and condition, let $K$ be the number of supported
requests and let $C$ be the number meeting the applicable completion
criterion. For $K>0$, the completion rate is $C/K$; matched conditions share $K$, so
their completion counts can also be compared directly. Unsupported requests
and their correct rejections are reported separately. Outcome distributions
over all arrivals retain both unsupported requests and failures. This
prevents successful rejections from inflating service completion.

\subsection{Latency and API Cost}
Decision latency is measured from client call initiation to receipt of the
full interpretation response. Admission waiting is accounted for separately
in the service timeline. We report medians and empirical 95th percentiles
(p95). The figures also show empirical cumulative distribution functions
(CDFs) and interquartile ranges (IQRs). These describe observed variation;
they are not confidence intervals.

For backend $m$ and request $i$, let $f_i^m$ be the time when its service
response reaches the originating controller. Its full request latency is
$T_i^m=f_i^m-a_i$, including admission waiting, interpretation, service
queueing, transfer, and execution. For a matched Jev--DeepSeek condition,
let $\mathcal{S}$ contain the requests completed correctly by both backends,
and let $N=|\mathcal{S}|$. The median paired saving
$\Delta_{\mathrm{paired}}$ is
\begin{equation}
\Delta_{\mathrm{paired}}=
\operatorname{median}_{i\in\mathcal{S}}
\left(T_i^{\mathrm{DeepSeek}}-T_i^{\mathrm{Jev}}\right).
\label{eq:paired}
\end{equation}
Positive values favor Jev. The percentage reduction of median latency uses
a different summary. Let $M_m$ be the median of $T_i^m$ over $\mathcal{S}$;
the reduction $R_T$, expressed as a percentage, is
\begin{equation}
R_T=100\left(1-\frac{M_{\mathrm{Jev}}}{M_{\mathrm{DeepSeek}}}\right).
\label{eq:latency-reduction}
\end{equation}
Equation~\ref{eq:paired} need not equal $M_{\mathrm{DeepSeek}}-M_{\mathrm{Jev}}$.
Both summaries are conditional on shared success, so we report $N$ and each
backend's completion count alongside them.

For a backend and condition, let $A$ be the provider-reported interpretation fees in
United States dollars (USD) and let $c$ be API cost per completion. Using
the corresponding completion count $C$, we define
\begin{equation}
c=\frac{A}{C},\qquad C>0.
\label{eq:cost}
\end{equation}
The numerator includes all interpretation calls for that condition, including
unsupported requests, unsuccessful jobs, and calls ending after their
deadline. It excludes warmups and separate semantic measurements. Study A
uses operational completion for $C$; Study B uses correct OCR completion.
Cost is undefined when $C=0$. Relative cost reductions use the same ratio
as Equation~\ref{eq:latency-reduction}, with $c$ replacing $M$.
These fees exclude computation, communication infrastructure, energy, and
maintenance.

\subsection{Analysis Scope}
The comparison units are matched conditions and consecutive measurement
blocks. Repeated texts, images, and cached interpretations make arrivals
dependent. We therefore report descriptive counts, quantiles, and differences.
Observed retention means matching or exceeding the measured completion count
under a matched condition. It does not establish statistical noninferiority:
no application-justified noninferiority margin was specified.
